% =============================================================================
% Template LaTeX IPG  —  https://github.com/VagnerBomJesus/TEMPLATE-IPG-LATEX
% Autor base, criador e mantenedor do projeto:
%     Vagner Bom Jesus  (@VagnerBomJesus)
%
% Proprietario e autor original deste template. Ao reutilizar, distribuir ou
% editar, MANTER este credito ao autor base (nao remover). Quem usa o template
% e o autor do SEU relatorio; o criador base do modelo e Vagner Bom Jesus.
% Licenca: GPL-3.0 (ver LICENSE).
% =============================================================================

% =============================================================================
% Capítulo 2 — Trabalho Relacionado
% Começar o capítulo numa nova página e incluir uma breve descrição sobre
% qual será o assunto do capítulo.
% Este capítulo deve apresentar o estudo realizado pelo autor na área do projeto.
% O autor deve demonstrar o quanto abrangente é o seu conhecimento e se o
% trabalho realizado foi desenvolvido com base nesse mesmo conhecimento.
% =============================================================================

\chapter{Trabalho Relacionado}
\label{cap:trabalho_relacionado}

Este capítulo apresenta o estado da arte relativo ao tema do trabalho, analisando os métodos, tecnologias e soluções existentes mais relevantes.

% ---- 2.1 Métodos e Tecnologias ----
\section{Métodos e Tecnologias Para \ldots}
\label{sec:metodos_tecnologias}

% Responder às seguintes questões:
% - Quais foram os métodos, ferramentas, tecnologias e conhecimento aplicado
%   usados para resolver o problema identificado?
% - O que já existe que possa ser reutilizado?
% Apresentar uma revisão da literatura relevante com referências bibliográficas.
% Usar \citep{chave} para citações entre parênteses ou \citet{chave} para
% citações no texto.
Descrever os métodos e tecnologias relevantes para o desenvolvimento do trabalho.

% ---- Exemplos de citação (substituir pelas referências reais) ----
%   \citet{chave}    -> "Autor (ano)"      (citação integrada no texto)
%   \citep{chave}    -> "(Autor, ano)"     (citação entre parênteses)
%   \textcite{chave} / \parencite{chave}   -> equivalentes nativos do biblatex
A título de exemplo, \citet{hillier2001} apresentam uma introdução abrangente à investigação operacional, enquanto o problema de dimensionamento de lotes foi formalizado por \citet{fleischmann1990}. Abordagens heurísticas para problemas relacionados têm sido amplamente estudadas na literatura \citep{pattloch2001, wiendhal1995}.

% ---- 2.2 Soluções Existentes ----
\section{Soluções Existentes Para \ldots}
\label{sec:solucoes_existentes}

% Responder às seguintes questões:
% - Qual foi a metodologia utilizada para recolher informações sobre soluções
%   existentes?
% - Quais são as soluções existentes?
% 